\documentclass[12pt, a4paper,journal, onecolumn, compsocar]{IEEEtran}
\IEEEoverridecommandlockouts
% The preceding line is only needed to identify funding in the first footnote. If that is unneeded, please comment it out.
\usepackage{cite}
\usepackage{amsmath,amssymb,amsfonts}
\usepackage{algorithmic}
\usepackage{graphicx}
\usepackage{textcomp}
\usepackage{xcolor}
\usepackage{fancyhdr}
\usepackage{mathtools}
\usepackage{lastpage}
\usepackage{hyperref}
\usepackage{caption}
\usepackage{cancel}
\captionsetup{justification=centering}
\usepackage{bm}
\usepackage{float}
\usepackage{thmtools}
\usepackage{thm-restate}
\hypersetup{
    colorlinks=true,
    linkcolor=blue,
    filecolor=magenta,      
    urlcolor=cyan,
    citecolor=blue,
}
\pagestyle{fancy}


\fancyhf{}
\fancyhead[C]{\small UNCuyo -- Ingeniería Mecatrónica / Mendoza -- Argentina. Espacio Curricular 311:\\ AUTOMÁTICA Y MÁQUINAS ELÉCTRICAS / PROYECTO GLOBAL INTEGRADOR.\\ Año: 2025. Alumnos: Calderón, Joaquín - Castel, Francisco. Fecha: 19 de Febrero 2026.}
\fancyfoot[C]{\small Pág. \thepage \hspace{1pt} de \pageref{LastPage}}
\renewcommand{\headrulewidth}{0.4pt}
\renewcommand{\footrulewidth}{0.4pt}

\setlength{\headheight}{45pt}
\def\BibTeX{{\rm B\kern-.05em{\sc i\kern-.025em b}\kern-.08em
    T\kern-.1667em\lower.7ex\hbox{E}\kern-.125emX}}
\makeatletter
\let\ps@plain\ps@fancy
\makeatother

% --- AJUSTE DE ESPACIADO DE TEXTO ---
\setlength{\parskip}{1em} % Añade espacio entre párrafos (y por ende, entre secciones)
\renewcommand{\baselinestretch}{1.15} % Aumenta un 15% el espacio entre líneas (más legible)

\begin{document}

% --- AJUSTE DE ESPACIADO DE ECUACIONES (GLOBAL) ---
\setlength{\abovedisplayskip}{12pt plus 3pt minus 3pt}      % Espacio antes de ecuación
\setlength{\belowdisplayskip}{12pt plus 3pt minus 3pt}      % Espacio después de ecuación
\setlength{\abovedisplayshortskip}{10pt plus 3pt minus 3pt} % Espacio corto antes
\setlength{\belowdisplayshortskip}{10pt plus 3pt minus 3pt} % Espacio corto después

\title{Proyecto Global Integrador: Control de Accionamiento de CA con
Motor Sincrónico de Imanes Permanentes}

\author{\IEEEauthorblockN{1\textsuperscript{st} Francisco Castel}
\IEEEauthorblockA{\textit{Facultad de Ingeniería} \\
\textit{UNCUYO}\\
Mendoza, Argentina \\
castel.francisco@uncuyo.edu.ar}
\and
\IEEEauthorblockN{2\textsuperscript{nd} Joaquín Calderón}
\IEEEauthorblockA{\textit{Facultad de Ingeniería} \\
\textit{UNCUYO}\\
Mendoza, Argentina \\
calderon.joaquin@uncuyo.edu.ar}
}

\maketitle
\thispagestyle{fancy}
\begin{abstract}
Modelado, simulación, diseño y análisis de desempeño de un Sistema de Control Automático de
Posición y Movimientos para un Accionamiento Eléctrico de 4 cuadrantes, compuesto por: máquina
eléctrica de corriente alterna (CA) trifásica sincrónica con excitación por imanes permanentes (PMSM),
alimentada por un inversor electrónico trifásico desde fuente de tensión continua (CC); reductor de
velocidad con engranajes planetarios conectado a la carga mecánica; retroalimentación con 1 sensor
de posición en el eje del motor más 3 sensores de corriente instantánea de fases en la salida del
inversor trifásico al estator de la máquina PMSM y 1 sensor de temperatura en el bobinado de estator.
\end{abstract}

\begin{IEEEkeywords}
accionamiento CA, simulación, térmico, cascada, control, modelado
\end{IEEEkeywords}
\section{Introducción}
El presente trabajo constituye el Proyecto Global Integrador del espacio curricular "Automática y Máquinas Eléctricas", cuyo objetivo principal es la integración de conocimientos y competencias fundamentales mediante su aplicación en un sistema mecatrónico.
El sistema bajo estudio consiste en un accionamiento eléctrico de cuatro cuadrantes compuesto por una máquina de corriente alterna trifásica sincrónica con excitación por imanes permanentes (PMSM), alimentada por un inversor electrónico. Este sistema acciona, a través de un reductor de velocidad, una carga mecánica correspondiente a una articulación de un brazo manipulador robótico elemental de un grado de libertad (péndulo rígido).
El alcance del proyecto abarca el modelado matemático, la simulación en Matlab/Simulink, y el diseño y análisis de desempeño de un sistema de control automático de posición y movimientos.

\section{Análisis, Modelado y Simulación del subsistema físico a Lazo Abierto}
Modelado de la carga mecánica, péndulo rígido de 1 g.d.l., fijo a una base en el sistema de referencia inercial, sometido a la fuerza de gravedad \(g\) con carga útil variable.

% 1. ECUACIONES DEL PÉNDULO
\begin{align}
    J_l\cdot \frac{d\omega}{dt} = T_q - b_l\cdot \omega_l - T_l(t)  \\
    T_l(t) = g\cdot k_l \cdot \sin(\theta_l(t)) + T_{ld}(t)
\end{align}

% 2. TEXTO EXPLICATIVO (Sin interrupciones)
Donde $\omega_l(t)$ y $\theta_l(t)$ representan la velocidad y posición angular de la articulación respectivamente; $T_q(t)$ es el torque impulsor entregado por el reductor y $T_{ld}(t)$ modela las perturbaciones externas variables. Los parámetros físicos constantes que caracterizan al péndulo y la transmisión se detallan con sus valores numéricos en la Tabla \ref{tab:parametros}.

% 3. CONTINUACIÓN LÓGICA (Inmediatamente después)
Es evidente que la relación cinemática viene dada por:
\[
    \frac{d\theta_l(t)}{dt} = \omega_l(t) \Leftrightarrow \theta_l(t) = \int_0^t{\omega_l(\xi)d\xi + \theta_l(0)}
\]

% 4. CÓDIGO DE LA TABLA (Al final, para que flote libremente)
\begin{table}[htbp]
    \caption{Parámetros Físicos y Especificaciones del Sistema}
    \centering
    \renewcommand{\arraystretch}{1.2} % Da un poco más de aire a las filas
    \begin{tabular}{|c|l|c|c|}
        \hline
        \textbf{Símbolo} & \textbf{Descripción} & \textbf{Valor} & \textbf{Unidad} \\
        \hline
        \multicolumn{4}{|c|}{\textbf{Subsistema Mecánico (Carga y Transmisión)}} \\
        \hline
        $r$ & Relación de transmisión & 120:1 & - \\
        $g$ & Aceleración de la gravedad & 9.81 & $\text{m}/\text{s}^2$ \\
        $m$ & Masa del brazo manipulador & 1.0 & kg \\
        $l_{cm}$ & Longitud al centro de masa & 0.25 & m \\
        $J_{cm}$ & Inercia centro de masa brazo & 0.0208 & $\text{kg}\cdot\text{m}^2$ \\
        $m_l$ & Masa de carga útil (variable) & $0 - 1.5$ & kg \\
        $l_{l}$ & Longitud total al extremo & 0.50 & m \\
        $b_{l}$ & Fricción viscosa articulación & 0.1 & $\text{Nms}/\text{rad}$ \\
        \hline
        \multicolumn{4}{|c|}{\textbf{Subsistema Mecánico (Motor)}} \\
        \hline
        $J_m$ & Inercia rotor (motor + caja) & $14 \times 10^{-6}$ & $\text{kg}\cdot\text{m}^2$ \\
        $b_m$ & Fricción viscosa (motor + caja) & $15 \times 10^{-6}$ & $\text{Nms}/\text{rad}$ \\
        $J_{eq}$ & Inercia equiv. total (Calc.) & $\approx 1.54 \times 10^{-5}$ & $\text{kg}\cdot\text{m}^2$ \\
        $b_{eq}$ & Fricción equiv. total (Calc.) & $\approx 2.19 \times 10^{-5}$ & $\text{Nms}/\text{rad}$ \\
        \hline
        \multicolumn{4}{|c|}{\textbf{Subsistema Eléctrico (PMSM)}} \\
        \hline
        $P_p$ & Pares de polos & 3 & - \\
        $\lambda_r$ & Flujo magnético de imanes & 0.016 & Wb \\
        $L_d$ & Inductancia eje directo & 6.6 & mH \\
        $L_q$ & Inductancia eje cuadratura & 5.8 & mH \\
        $L_{ls}$ & Inductancia de dispersión & 0.8 & mH \\
        $R_{s}$ & Resistencia estator (@ $20^\circ$C) & 1.02 & $\Omega$ \\
        \hline
        \multicolumn{4}{|c|}{\textbf{Subsistema Térmico}} \\
        \hline
        $C_{ts}$ & Capacitancia térmica estator & 0.818 & $\text{J}/^\circ\text{C}$ \\
        $R_{th}$ & Resitencia térmica Estator-Amb & 146.7 & $^\circ\text{C}/\text{W}$ \\
        $\alpha_{cu}$ & Coef. temp. del Cobre & $3.9 \times 10^{-3}$ & $1/^\circ\text{C}$ \\
        $T_{amb}$ & Temperatura ambiente & 20 & $^\circ\text{C}$ \\
        \hline
        \multicolumn{4}{|c|}{\textbf{Límites de Operación (Especificaciones)}} \\
        \hline
        $n_{l,nom}$ & Velocidad nominal de carga & 60 & rpm \\
        $T_{q,nom}$ & Torque nominal de salida & 17.0 & Nm \\
        $T_{q,max}$ & Torque máximo (pico) & 45.0 & Nm \\
        \hline
    \end{tabular}
    \label{tab:parametros}
\end{table}

\subsection{Tren de transmisión}
Como se explicó más arriba, la velocidad del motor no tendrá relación 1:1 con el péndulo, por lo que se implementa una \textbf{tren de transmisión planetario} como caja reductora reversible. Se asume lo siguiente:
\begin{itemize}
    \item Ausencia de \textit{backlash}.
    \item Ausencia de elasticidad torcional (acoplamiento rígido).
    \item Momento de inercia equivalente (parámetros concentrados).
    \item Pérdidas de energía por fricción interna consideradas en el eje de entrada junto al motor.
\end{itemize}
Por lo tanto, debido a las simplificaciones realizadas, es válido considerar a la carga mecánica y el motor eléctrico como una unidad en sí misma, lo que permite tener un subsistema mecánico de un grado de libertad, permitiendo escribir las siguientes ecuaciones.
\begin{align}
    \omega_l(t) = \frac{1}{r}\omega_m(t) \\
    T_q(t) = rT_d(t)
\end{align}
tal que:
\begin{itemize}
    \item r: Relación de velocidades del tren de transmisión.
    \item \(\omega_m(t)\): Velocidad angular del eje del motor.
    \item \(T_d(t)\): Torque de entrada al tren de transmisión.
    \item \(T_q(t)\): Torque de salida al tren de transmisión.
    
\end{itemize}

    \begin{figure}[htbp]
        \centering
        \includegraphics[width=0.6\linewidth]{imagenes/brazo_1gdl.png}
        \caption{Robot manipulador elemental de 1 g.d.l en plano vertical (péndulo rígido actuado)}
        \label{fig:placeholder}
    \end{figure}

\subsection{Máquina PMSM}
    Para el accionamiento, se utilizará una máquina eléctrica de CA trifásica sincrónica con excitación por imanes permanentes (\textbf{PMSM}) y estator simétrico y equilibrado, con bornes de fases \textit{abc} (coordenadas reales y "estacionarias) conectados en configuración estrella (Y) con neutro flotante \textbf{no} accesible.
   \subsubsection{Modelo del subsistema mecánico del motor}
   El modelo matemático equivalente del subsistema mecánico del rotor de la máquina eléctrica (referido a estator estacionario \(\equiv\) fijo al sistema inercial de referencia) se describe mediante las siguientes \textbf{ecuaciones}.
    \begin{equation}\label{EqJm}
        J_m \frac{d\omega_m}{dt} = T_m - b_m\omega_m - T_d 
    \end{equation}
    \begin{equation}
        \frac{d\theta_m(t)}{dt} = \omega_m(t) \Leftrightarrow \theta_m(t) = \int_0^t{\omega_m(\xi)d\xi} + \theta_m(0)
    \end{equation}
    \subsubsection{Modelo del subsistema electromagnético del motor}
    Modelo idealizado equivalente en coordenadas electromagnéticas de rotor \(d\theta_r\) "fijas" al rotor, aplicando la \textit{Transformación de Park} al circuito de estator estacionario (marco de referencia virtual eléctrico de rotor, sincrónico solo en régimen estacionario).
    \begin{gather}
        \frac{d\theta_r(t)}{dt} \equiv \omega_r(t) = \theta_r(t) = \int_0^t{\omega_r(\xi)d\xi+\theta_r(0)} \\
        \theta_r(t) \equiv P_p\theta_m(t) \therefore \omega_r(t) = P_p\omega_m(t)
    \end{gather}
    \textbf{Torque electromagnético} = \textit{torque debido a imanes + torque de reluctancia}:
    \begin{restatable}{equation}{torqueElectromagnetico}
            T_m(t) = \frac{3}{2} P_p \cdot \lambda'^r_m \cdot i^r_{qs}(t) + \frac{3}{2} P_p (L_d-L_q)\cdot i^r_{ds}(t)\cdot i^r_{qs}(t)
            \label{eq:torque_electromagnetico}
    \end{restatable}
    % \begin{equation}\label{eq:torque_electromagnetico}
    % T_m(t) = \frac{3}{2} P_p \cdot \lambda'^r_m \cdot i^r_{qs}(t) + \frac{3}{2} P_p (L_d-L_q)\cdot i^r_{ds}(t)\cdot i^r_{qs}(t)
    % \end{equation}
 
    \textbf{Ecuaciones de tensiones de fase referidas al rotor}
    \begin{align}
        K_s\textbf{v}_{abcs} = \textbf{v}^r_{qd0s}
    \end{align}
    Siendo \(K_s\) la matriz de transformación de Park.
    \begin{gather}
        v^r_{qs}=R_s(T_s^\circ(t))i^r_{qs}(t) + L_q\frac{d i^r_{qs}(t)}{dt}+[\lambda'^r_m + L_d i^r_{ds}(t)]\omega_r(t)\\
        v_{ds}^r=R_s(T_s^\circ(t))i^r_{ds}(t)+L_d\frac{d i^r_{ds}(t)}{dt} - L_qi^r_{qs}(t)\omega_r(t)\\ 
        v^r_{0s}=R_s(T_s^\circ(t))i^r_{0s}(t)+L_{ls}\frac{d i^r_{0s}(t)}{dt}
    \end{gather}
    \subsubsection{Modelo del subsistema térmico}
    Considerando también la influencia de la temperatura en la resistencia de los devanados se explicitan las ecuaciones que la rigen considerando solo pérdidas resistivas por efecto Joule y transferencia de calor solo por conducción. \\
    \textbf{Potencia de pérdidas en sistema de ref. rotacional}
    \begin{gather}
        P_{s \ perd}(t) = R_s(T_s^\circ(t)) \frac{3}{2}\cdot\Big({i^r_{qs}}^2(t) +{i^r_{ds}}^2(t) + 2{i_{0s}}^2(t)\Big) 
    \end{gather}
    \textbf{Balance térmico del estator}
    \begin{gather}
        P_{s \ perd}(t) = C_{ts}\frac{d T_s^\circ(t)}{dt} + \frac{1}{R_{ts-amb}}\cdot\Big(T_s^\circ(t) - T^\circ_{amb}(t)\Big)
    \end{gather}
    
    \subsection{Modelo matematico equivalente del subsistema mecánico}
    Habiendo expresado las ecuaciones para los distintos subsistemas de la máquina PMSM y del tren de transmisión, es posible conjugarlas en un modelo equivalente.
    Sustituyendo \(\omega_l(t)\) de (3) y \(T_q(t) = rT_d(t)\) de (4) en la Ec. (1) tenemos:
    \begin{equation}\label{EqAux1}
        \frac{J_l}{r}\frac{d\omega_n(t)}{dt} = rT_d(t)-\frac{b_l}{r}\omega_m(t)-T_l(t)
    \end{equation}
    Despejando \(T_d(t)\) y reemplazando \(T_l(t)\) por la Eq. 2:
    \begin{equation}\label{Eq1}
        T_d(t)=\frac{J_l}{r^2}\frac{d\omega_m}{dt}+\frac{b_l}{r^2}\omega_m+\frac{1}{r}\Big(g\cdot k_l\cdot \sin{\Big(\frac{\theta_m(t)}{r}\Big)+T_{ld}(t)}\Big)
    \end{equation}
    Reemplazando \(T_d(t)\) de \eqref{Eq1} en  \eqref{EqJm}:
    \begin{equation}
        J_{eq}\frac{d\omega_m}{dt}=T_m-b_m\omega_m-\frac{1}{r}\Big(g\cdot k_l \sin{\Big(\frac{\theta_m(t)}{r}\Big)} + T_{ld}(t)\Big)
    \end{equation}
    \begin{gather}
        J_{eq} = J_m +\frac{J_l}{r^2}\\
        b_{eq} = b_m + \frac{b_l}{r^2}\\
        \theta_l(t)=\frac{\theta_m(t)}{r}
    \end{gather}
    Podemos entonces, a partir de lo desarrollado, expresar el modelo en el espacio de estados.\\
    \textbf{Ecuaciones de estado:} Las ecuaciones de estado del subsistema mecánico completo quedan:
    \begin{equation}
        \begin{cases}
            \dot \theta_m(t) = \omega_m(t), \\ 
            \dot\omega_m(t)=-\frac{b_{eq}}{J_{eq}}\omega_m(t)-\frac{g\cdot k_l}{J_{eq}r}\cdot \sin{\Big(\frac{\theta_m(t)}{r}\Big)}+\frac{1}{J_{eq}}T_m(t)-\frac{1}{J_{eq}r}T_{ld}(t),\\
            y(t)=\theta_m(t)
        \end{cases}
        \label{eq: EqEstadoModeloMecanico}
    \end{equation}
    \textbf{Forma matricial}
    \begin{equation}
        \begin{cases}
            \begin{bmatrix}
                \dot\theta_m(t) \\\dot\omega_m(t)
            \end{bmatrix}
            = 
            \begin{bmatrix}
                \omega_m(t)\\
                -\frac{g\cdot k_l}{J_{eq}r}\cdot\sin{\big(\frac{\theta_m(t)}{r}\big)}-\frac{b_{eq}}{J_{eq}}\cdot\omega_m(t)
            \end{bmatrix}
            +
            \begin{bmatrix}
                0 & 0\\
                \frac{1}{J_{eq}} &-\frac{1}{J_{eq}r}
            \end{bmatrix}
            \begin{bmatrix}
                T_m(t)\\
                T_{ld}(t)
            \end{bmatrix}\\
            y(t) = \begin{bmatrix}
                1 & 0
            \end{bmatrix}
            \begin{bmatrix}
                \theta_m(t)\\
                \omega_m(t)
            \end{bmatrix}
        \end{cases}
    \end{equation}
    \subsection{Modelo dinámico del sistema físico completo}
    Asumiendo, para el caso general, \(i_{0s}^r \neq 0\):\\
    
    \textbf{Variables de estado:} \[\bm{x} =[i_{qs}^r(t),i_{ds}^r(t),i_{0s}^r(t), \omega_m(t),\theta_m(t),T_s^\circ]^T\] 
    Despejando de (11), (12) y (13) obtenemos: 
    \begin{equation}
    \begin{aligned}
        \frac{di^r_{qs}}{dt} &= \frac{v^r_{qs}}{L_q}-\frac{1}{L_q}R_s(T_s^\circ(t))i^r_{qs}-\frac{1}{L_q}\Big([\lambda'^r_{m}+L_di^r_{ds}]\omega_r(t)\Big) \\
        \frac{di^r_{ds}}{dt} &= \frac{v^r_{ds}}{L_d}-\frac{1}{L_d}R_s(T_s^\circ(t))i^r_{ds}+\frac{L_q}{L_d}i^r_{qs}\omega_r(t) \\
        \frac{di^r_{0s}}{dt} &= \frac{v^r_{0s}}{L_{ls}}-\frac{1}{L_{ls}}R_s(T_s^\circ(t))i^r_{0s}(t)
    \end{aligned}
    \end{equation}
    Luego, a partir del \textit{balance térmico del estator} (15) y de la \textit{potencia de pérdida} en el sistema de coordenadas rotativo tenemos: 
    \begin{equation}
        \begin{aligned}
        R_s(T_s^\circ(t)) = R_{sREF} \big(1+\alpha_{Cu}(T_s^\circ(t)-T^\circ_{sREF})\big)\\
            P_{s\ perd}(t) = R_s(T_s^\circ(t))\frac{3}{2}\Big({i^r_{qs}(t)}^2+{i^r_{ds}(t)}^2+2\cdot{i^r_{0s}(t)}^2\Big)  \\
            \frac{dT_s^\circ(t)}{dt}=\frac{P_{s \ perd}(t)}{C_{ts}}-\frac{1}{R_{ts-amb}\cdot C_{ts}}\Big(T_s^\circ(t)-T_{amb}^\circ(t)\Big)
        \end{aligned}
    \end{equation}

    \subsection{Resumen de Ecuaciones de Estado del Sistema Completo}
    
    A modo de síntesis, el comportamiento dinámico total del accionamiento (Mecánico + Eléctrico + Térmico) queda definido por el siguiente sistema de 6 ecuaciones diferenciales no lineales de primer orden. Este conjunto de ecuaciones constituye la base matemática para la linealización y el posterior diseño del control.
    \begin{equation}
    \left\{
    \begin{aligned}\label{eq:equaciones de estado}
           \frac{d\theta_m}{dt} &= \omega_m \\
        \frac{d\omega_m}{dt} &= \frac{1}{J_{eq}} \left[ \frac{3}{2} P_p \lambda_{pm} i_{qs}^r -
        B_{eq}\omega_m - \frac{g k_l}{r} 
        \sin\left(\frac{\theta_m}{r}\right) - T_{ld} \right]   \\
        \frac{di^r_{qs}}{dt} &= \frac{1}{L_ q}\Big[v^r_{qs}-R_s(T_s^\circ(t))i^r_{qs}-[\lambda'^r_{m}+L_di^r_{ds}]\omega_r(t)\Big] \\
        \frac{di^r_{ds}}{dt} &= \frac{1}{L_d}\Big[v^r_{ds}-R_s(T_s^\circ(t))i^r_{ds}+L_qi^r_{qs}\omega_r(t)\Big] \\
        \frac{di^r_{0s}}{dt} &= \frac{1}{L_{ls}}\Big[v^r_{0s}-R_s(T_s^\circ(t))i^r_{0s}(t)\Big]    \\
        \frac{dT_s^\circ}{dt} &= \frac{1}{C_{ts}} \left[ \frac{3}{2} R_s(T_s) \left( (i_{qs}^r)^2 + (i_{ds}^r)^2 \right) - \frac{T_s^\circ - T_{amb}^\circ}{R_{th}} \right]
    \end{aligned}
    \right.
    \end{equation}
    Donde las variables de estado son $\mathbf{x} = [i_{qs}^r, i_{ds}^r, i_{0s}^r, \omega_m, \theta_m, T_s^\circ]^T$ y las entradas de control son $\mathbf{u} = [v_{qs}^r, v_{ds}^r, v_{0s}^r]^T$.

% ==========================================================================
% SECCIÓN II-E: VALIDACIÓN POR SIMULACIÓN (LAZO ABIERTO)
% ==========================================================================

\subsection{Implementación y Validación en Entorno de Simulación}
Se implementó el modelo matemático completo en el entorno MATLAB/Simulink, estructurando el sistema mediante bloques funcionales interconectados que representan las ecuaciones diferenciales deducidas previamente.

La arquitectura de la simulación se diseñó para validar tanto la dinámica de la planta como la coherencia de las transformaciones de coordenadas. El flujo de señales se configura de la siguiente manera:
\begin{enumerate}
    \item \textbf{Generación de Consigna:} Se establecen tensiones de referencia constantes en el marco rotórico ($v_q^*=5$ V, $v_d^*=0$).
    \item \textbf{Transformaciones de Coordenadas:} Las tensiones $dq$ se convierten al marco trifásico $abc$ mediante la \textit{Transformada Inversa de Park} y se reconvierten inmediatamente al marco $dq$ mediante la \textit{Transformada de Park}. Esta estructura en bucle ("Back-to-Back") permite verificar la integridad matemática de los algoritmos antes de ingresar a la planta.
    \item \textbf{Planta Física:} Las señales resultantes alimentan al \textbf{Subsistema Electromecánico} y al \textbf{Subsistema Mecánico}, interactuando simultáneamente con el \textbf{Subsistema Térmico}.
\end{enumerate}

% FIGURA DIAGRAMA SIMULINK
\begin{figure}[H]
    \centering
    \includegraphics[width=1.0\linewidth]{imagenes/simulink_diagrama_completo.png}
    \caption{Diagrama de bloques implementado en Simulink. Se destaca la estructura de validación con transformadas directa e inversa y la interconexión de los subsistemas.}
    \label{fig:simulink_top}
\end{figure}

\subsubsection{Resultados del Ensayo a Lazo Abierto}
Se realizó un ensayo de excitación al escalón con $v_q = 5$ V y carga nula ($T_{ld}=0$), partiendo de condiciones iniciales de reposo y equilibrio térmico ($T_s(0)=20^\circ$C).

\textbf{Análisis Electromecánico:}
La Fig. \ref{fig:res_elec_mec} vincula la dinámica eléctrica con la mecánica:
\begin{itemize}
    \item \textbf{Superior ($i_q$):} Al aplicar la tensión, la corriente aumenta rápidamente ($\tau_e = L_q/R_s$). Sin embargo, a medida que el motor gana velocidad, la aparición de la fuerza contraelectromotriz ($e = \lambda_r \omega_r$) se opone a la tensión aplicada, provocando la reducción de la corriente observada a partir de $t \approx 0.1$ s.
    \item \textbf{Inferior ($\omega_m$):} El torque generado por la corriente acelera el rotor venciendo la inercia $J_{eq}$. La velocidad se estabiliza cuando el torque motor se iguala con el torque de fricción viscosa.
\end{itemize}

% FIGURA COMBINADA CORRIENTE/VELOCIDAD (CLAVADA AQUÍ CON [H])
\begin{figure}[H]
    \centering
    \includegraphics[width=0.9\linewidth]{imagenes/resultado_corriente.png} \\[0.1cm]
    \includegraphics[width=0.9\linewidth]{imagenes/resultado_velocidad.png}
    \caption{Respuesta dinámica al escalón de tensión $v_q=5$V. \textbf{Arriba:} Corriente de estator $i_q$ [A], mostrando el transitorio inductivo y el efecto de la Back-EMF. \textbf{Abajo:} Velocidad mecánica del rotor $\omega_m$ [rad/s].}
    \label{fig:res_elec_mec}
\end{figure}

\textbf{Análisis Térmico:}
La Fig. \ref{fig:res_term} muestra la evolución de la temperatura del estator. Se observa un inicio correcto en $T_{amb}=20^\circ$C.
Es notable que la pendiente de calentamiento no es constante: es máxima al inicio (cuando la corriente es alta y el rotor está detenido) y disminuye cuando el motor acelera y la corriente baja debido a la fuerza contraelectromotriz. Este comportamiento valida el correcto acoplamiento entre las variables eléctricas ($i^2 R$) y la dinámica mecánica en el modelo térmico.

% FIGURA TÉRMICA (CLAVADA AQUÍ CON [H])
\begin{figure}[H]
    \centering
    \includegraphics[width=0.9\linewidth]{imagenes/resultado_temperatura.png}
    \caption{Evolución de la temperatura del estator $T_s$. Nótese el cambio de pendiente provocado por la variación de la corriente durante el arranque.}
    \label{fig:res_term}
\end{figure}

% ESTE COMANDO ASEGURA QUE LAS FOTOS TERMINEN AQUÍ Y NO PASEN A LA MATEMÁTICA
\clearpage

% ==========================================================================
% SECCIÓN II-F: DEFINICIÓN LPV
% ==========================================================================

\subsection{Modelo global linealizado con parametros variables (LPV)}
A partir del modelo no lineal (NL) validado, se aplica la Linealización Jacobiana mediante aproximación con Serie de Taylor truncada en 1er orden.
Considerando que el actual modelo es de la forma $\bm{\dot x} =\bm{f}(\bm{x},u)$, se define el sistema lineal incremental respecto a un punto de equilibrio ($\bm{x}_0, \bm{u}_0$):
\begin{equation*}
    \delta\dot x=A_0\delta \bm{x}+ B_0\delta u
\end{equation*}

Las matrices de entrada manipulada y de perturbación, al ser independientes de los estados en esta estructura, resultan constantes en todo el espacio de operación:
\begin{equation*}
B_0^c = \begin{bmatrix}
0 & 0 & 0 \\
0 & 0 & 0 \\
\frac{1}{L_q} & 0 & 0 \\
0 & \frac{1}{L_d} & 0 \\
0 & 0 & \frac{1}{L_{ls}} \\
0 & 0 & 0
\end{bmatrix} , \quad
B_0^d = \begin{bmatrix}
0 & 0 \\
-\frac{1}{J_{eq} r} & 0 \\
0 & 0 \\
0 & 0 \\
0 & 0 \\
0 & \frac{1}{C_{ts} R_{th}}
\end{bmatrix} 
\end{equation*}

La matriz del sistema $\mathbf{A}_0$, que contiene la dinámica y los acoplamientos electromecánicos y térmicos variables, se analiza en detalle a continuación.

%-------------- LPV -----------------
\subsection{Modelo global linealizado con parametros variables (LPV) con \(i^r_{ds}\neq0\)}
A partir del modelo NL anteriormente presentado, es posible a través de la Linealización Jacobiana, mediante aproximación con Serie de Taylor truncada en 1er orden en un punto genérico de operación, obtener un modelo lineal mas simple de analizar y controlar.

Considerando que el actual modelo es de la forma \[\bm{\dot x} =\bm{f}(\bm{x},u)\]
Donde lo primero es encontrar los puntos de equilibrio, del sistema, es decir los valores \(\bm{x}_0 ,\)\(u_0 ,\) tal que \(\bm{\dot x_0} = \bm{f}(\bm{x}_0,u_0)\), y definimos \(x=x_0 +\delta x\) y \(u=u_0+\delta u\). Luego, expandimos la ecuación no lineal en términos de las perturbaciones respecto a estos valores de equilibrio, lo que resulta en: 
\begin{equation*}
    \dot x_0 + \delta \dot x \cong \bm{f}(\bm{x}_0,u_0) + A \delta x+B\delta u
\end{equation*}
Donde \(A\) y \(B\) son la mejor aproximación lineal de la función \(\bm{f}(x,u) \) en \(x_0, u_0\), calculándose: 
\begin{equation*}
    \bm{A}_0 = \left[ \frac{\partial \bm{f}}{\partial \bm{x}} \right]_{\bm{x}_o, u_o} \quad \text{y} \quad \bm{B}_0 = \left[ \frac{\partial \bm{f}}{\partial u} \right]_{\bm{x}_o, u_o}
\end{equation*}
Restando la solución de equilibrio y dejando la ecuación solo en variables incrementales nos queda:
\begin{equation*}
    \delta\dot x=A_0\delta \bm{x}+ B_0\delta u
\end{equation*}
Es decir, una ecuación diferencial lineal alrededor del punto de equilibrio cuasi-estático.
\subsubsection{Implementación en caso genérico de operación, \(\omega=cte\)}
Utilizando entonces las ecuaciones vistas podemos comenzar a linealizar nuestro sistema mediante el método de pequeñas señales.
Primero, encontramos nuestro espacio de operación global no lineal, que es el espacio donde se resolvera la primer parte de la ecuación expandida, \(\bm{f}(\bm{x}_0,u_0)=0\)
\begin{equation}
\left\{
    \begin{aligned}
        \frac{d\theta_{mO}}{dt}(t) &= \omega_{mO}(t) = \text{cte} \\[10pt]
        %
        \frac{d\omega_{mO}}{dt}(t) &= \frac{1}{J_{eq}} \left[ -b_{eq}\omega_{mO}(t) + \frac{3}{2}P_p \left[ \lambda_m' + (L_d - L_q)i_{dsO}^r(t) \right] i_{qsO}^r(t) - g \cdot k_l \sin\left(\frac{\theta_{mO}}{r}\right) - \frac{1}{r}T_{ld}(t) \right] = 0 \\[10pt]
        %
        \frac{di_{qsO}^r}{dt}(t) &= \frac{1}{L_q} \left[ v_{qsO}^r(t) - R_s(T_{sO}^\circ(t))i_{qsO}^r(t) - (\lambda_m' + L_d i_{dsO}^r(t))P_p \omega_{mO}(t) \right] = 0 \\[10pt]
        %
        \frac{di_{dsO}^r}{dt}(t) &= \frac{1}{L_d} \left[ v_{dsO}^r(t) - R_s(T_{sO}^\circ(t))i_{dsO}^r(t) + L_q i_{qsO}^r(t)P_p \omega_{mO}(t) \right] = 0 \\[10pt]
        %
        \frac{di_{0sO}^r}{dt}(t) &= \frac{1}{L_{ls}} \left( v_{0sO}^r(t) - R_s(T_{sO}^\circ(t))i_{0sO}^r(t) \right) = 0 \\[10pt]
        %
        \frac{dT_{sO}^\circ}{dt}(t) &= \frac{1}{C_{ts}} \left[ \frac{3}{2}R_s(t) \left( i_{qsO}^r(t)^2 + i_{dsO}^r(t)^2 + 2i_{0sO}^r(t)^2 \right) - \frac{1}{R_{ts-amb}} (T_{sO}^\circ(t) - T_{ambO}^\circ(t)) \right] = 0
    \end{aligned}
\right.
\end{equation}
Luego la parte lineal dinámica, que es de principal interés en este análisis.
\begin{equation*}
    \Delta \dot{x}(t) = \underbrace{\left. \frac{\partial f}{\partial x} \right|_0}_{A_o} (t) \cdot \Delta x(t) + \underbrace{\left. \frac{\partial f}{\partial u} \right|_0}_{B_o} (t) \cdot \Delta u(t); \quad \Delta x(0) = 0
\end{equation*}
\begin{equation*}
\left\{
    \begin{aligned}
        \Delta \dot{\theta}_m(t) &= \Delta \omega_m(t) \\[10pt]
        %
        \Delta \dot{\omega}_m(t) &= \frac{1}{J_{eq}} \left[ \frac{3}{2} P_p \left\{ [\lambda_m' + (L_d - L_q) \cdot i_{dsO}^r] \cdot \Delta i_{qs}^r(t) + (L_d - L_q) \cdot i_{qsO}^r \cdot \Delta i_{ds}^r(t) \right\} - b_{eq} \cdot \Delta \omega_m(t) \right. \\
        &\quad \left. - \frac{g \cdot k_l}{r^2} \cdot \cos\left(\frac{\theta_{mO}}{r}\right) \cdot \Delta \theta_m(t) - \frac{\Delta T_{ld}(t)}{r} \right] \\[10pt]
        %
        \Delta \dot{i}_{qs}^r(t) &= \frac{1}{L_q} \left[ \Delta v_{qs}^r(t) - R_{sO} \cdot \Delta i_{qs}^r(t) - R_{s,REF} \alpha_{cu} \Delta T_s^\circ(t) \cdot i_{qsO}^r - \{\lambda_m' + L_d \cdot i_{dsO}^r\} P_p \cdot \Delta \omega_m(t) \right. \\
        &\quad \left. - L_d P_p \omega_{mO} \cdot \Delta i_{ds}^r(t) \right] \\[10pt]
        %
        \Delta \dot{i}_{ds}^r(t) &= \frac{1}{L_d} \left[ \Delta v_{ds}^r(t) - R_{sO} \cdot \Delta i_{ds}^r(t) - R_{s,REF} \alpha_{cu} \Delta T_s^\circ(t) \cdot i_{dsO}^r + L_q \cdot i_{qsO}^r \cdot P_p \cdot \Delta \omega_m(t) \right. \\
        &\quad \left. + L_q \cdot P_p \cdot \omega_{mO} \cdot \Delta i_{qs}^r(t) \right] \\[10pt]
        %
        \Delta \dot{i}_{0s}^r(t) &= \frac{1}{L_{ls}} \left[ \Delta v_{0s}^r(t) - R_{sO} \cdot \Delta i_{0s}^r(t) - R_{s,REF} \alpha_{cu} \Delta T_s^\circ(t) \cdot i_{0s}^r \right] \\[10pt]
        %
        \Delta \dot{T}_s(t) &= \frac{1}{C_{ts}} \left[ \frac{3}{2} \cdot R_{sO} \cdot (2 \cdot i_{qsO}^r \cdot \Delta i_{qs}^r(t) + 2 \cdot i_{dsO}^r \cdot \Delta i_{ds}^r(t) + 4 \cdot i_{0sO}^r \cdot \Delta i_{0s}^r(t)) \right. \\
        &\quad + \frac{3 \alpha_{cu} R_{sREF}}{2} (i_{qsO}^r(t)^2 + i_{dsO}^r(t)^2 + 2i_{0sO}^r(t)^2) \cdot \Delta T_s(t) \\
        &\quad \left. - \frac{1}{R_{ts-amb}} (\Delta T_s(t) - \Delta T_{amb}(t)) \right]
    \end{aligned}
\right.
\end{equation*}
Donde \(R_{sO}=R_s(T^\circ_{sO})\)
\begin{equation*}
A_0 = \begin{bmatrix}
0 & 1 & 0 & 0 & 0 & 0 \\
-\frac{g k_l \cos\left(\frac{\theta_{mO}}{r}\right)}{J_{eq} r^2} & -\frac{b_{eq}}{J_{eq}} & \frac{3 P_p (\lambda_m + i_{dsO} (L_d - L_q))}{2 J_{eq}} & \frac{3 P_p i_{qsO} (L_d - L_q)}{2 J_{eq}} & 0 & 0 \\
0 & -\frac{P_p (\lambda_m + L_d i_{dsO})}{L_q} & -\frac{R_s}{L_q} & -\frac{L_d P_p \omega_{mO}}{L_q} & 0 & -\frac{R_{s,ref} \alpha_{cu} i_{qsO}}{L_q} \\
0 & \frac{L_q P_p i_{qsO}}{L_d} & \frac{L_q P_p \omega_{mO}}{L_d} & -\frac{R_s}{L_d} & 0 & -\frac{R_{s,ref} \alpha_{cu} i_{dsO}}{L_d} \\
0 & 0 & 0 & 0 & -\frac{R_s}{L_{ls}} & -\frac{R_{s,ref} \alpha_{cu} i_{0sO}}{L_{ls}} \\
0 & 0 & \frac{3 R_s i_{qsO}}{C_{ts}} & \frac{3 R_s i_{dsO}}{C_{ts}} & \frac{6 R_s i_{0sO}}{C_{ts}} & \Gamma_{O}
\end{bmatrix}
\end{equation*}

Donde la constante de disipación térmica acoplada $\Gamma_{th}$ se define como:
\begin{equation*}
\Gamma_{th} = -\frac{\frac{1}{R_{ts,amb}} - \frac{3 R_{s,ref} \alpha_{cu} (2 i_{0sO}^2 + i_{dsO}^2 + i_{qsO}^2)}{2}}{C_{ts}}
\end{equation*}
Las matrices de entrada manipulada y de perturbación quedan respectivamente:
\begin{equation*}
B_o^c = \begin{bmatrix}
0 & 0 & 0 \\
0 & 0 & 0 \\
\frac{1}{L_q} & 0 & 0 \\
0 & \frac{1}{L_d} & 0 \\
0 & 0 & \frac{1}{L_{ls}} \\
0 & 0 & 0
\end{bmatrix} ,
\end{equation*}

\begin{equation*}
B_o^d = \begin{bmatrix}
0 & 0 \\
-\frac{1}{J_{eq} r} & 0 \\
0 & 0 \\
0 & 0 \\
0 & 0 \\
0 & \frac{1}{C_{ts} R_{ts,amb}}
\end{bmatrix} 
\end{equation*}
El modelo global LPV toma la forma:
\begin{equation}
\begin{bmatrix}
\Delta \dot{\theta}_m(t) \\
\Delta \dot{\omega}_m(t) \\
\Delta \dot{i}_{qs}^r(t) \\
\Delta \dot{i}_{ds}^r(t) \\
\Delta \dot{i}_{0s}(t) \\
\Delta \dot{T}_s(t)
\end{bmatrix} = A_o \cdot 
\begin{bmatrix}
\Delta \theta_m(t) \\
\Delta \omega_m(t) \\
\Delta i_{qs}^r(t) \\
\Delta i_{ds}^r(t) \\
\Delta i_{0s}(t) \\
\Delta T_s(t)
\end{bmatrix} + B_o^c \cdot 
\begin{bmatrix}
\Delta v_{qs}^r(t) \\
\Delta v_{ds}^r(t) \\
\Delta v_{0s}(t)
\end{bmatrix} + B_o^d \cdot 
\begin{bmatrix}
\Delta T_{ld}(t) \\
\Delta T_{amb}(t)
\end{bmatrix} .
\end{equation}


\subsection{Linealización por Retroalimentación NL: Modelo LTI equivalente}
El objetivo de esta sección es obtener un modelo Lineal Invariante en el Tiempo (LTI) a partir de la planta no lineal, aplicando estrategias de Control Vectorial con campo orientado. Para lograr esto, se imponen restricciones matemáticas en las entradas manipuladas mediante retroalimentación de estado, buscando desacoplar los canales de flujo magnético y torque. Para el modelo LTI, se asume la dinámica térmica como lineal ($R_s = R_{s0}$ constante). Además, la no linealidad gravitacional de la carga se separa explícitamente de la matriz de estado, tratándola como una perturbación externa a ser compensada.

\subsubsection{Restricción o Ley de Control Mínima en el eje $d$}
Para lograr el desacoplamiento de flujo y torque, se impone el requisito de mantener la corriente del eje directo nula: $i_{ds}^r(t) \equiv 0$. Analizando la ecuación diferencial de dicho estado:
\begin{equation}
    \frac{di_{ds}^r(t)}{dt} = \frac{1}{L_d} \left[ v_{ds}^r(t) - R_{s0} i_{ds}^r(t) + L_q i_{qs}^r(t) \omega_r(t) \right]
\end{equation}

La \textbf{Ley de Control Mínima} requerida sobre la variable manipulada virtual $v_{ds}^r(t)$ para cancelar la fuerza contraelectromotriz cruzada y forzar la dinámica a cero (asumiendo la hipótesis de que el estado inicial es nulo, $i_{ds}^r(0) = 0$) es:
\begin{equation} \label{eq:vds_minima}
    v_{ds}^r(t) = -L_q \cdot i_{qs}^r(t) \cdot \omega_r(t)
\end{equation}

Al aplicar este voltaje, la dinámica del eje directo se anula y la ecuación de torque electromagnético se vuelve lineal y proporcional a $i_{qs}^r$.

\subsubsection{Ecuaciones Escalares y Matriciales del Modelo LTI (Nominal)}
Aplicando las leyes de control obtenidas \eqref{eq:vds_minima} y \eqref{eq:vqs_comp}, y asumiendo $i_{ds}^r(t) \equiv 0$, el sistema electromecánico global se reduce a un modelo LTI equivalente. El sistema de ecuaciones diferenciales resultante es:
\begin{equation} \label{eq:sis_escalar_LTI}
\begin{cases}
    \vspace{0.2cm}
    \dfrac{d\theta_m(t)}{dt} = \omega_m(t) \\
    \vspace{0.2cm}
    \dfrac{d\omega_m(t)}{dt} = \dfrac{1}{J_{eq}}\Big[\frac{3 P_p \lambda'_m}{2} i_{qs}(t) - b_{eq} \omega_m(t) - \frac{1}{r} T_{ld}(t) - \frac{g\cdot k_l}{r} \sin\left(\frac{\theta_m(t)}{r}\right)\Big] \\
    \dfrac{di_{qs}(t)}{dt} = \dfrac{1}{L_q}\Big[-P_p \lambda'_m \omega_m(t) - R_{sREF}\cdot i_{qs}(t) + v_{qs}(t)\Big]
\end{cases}
\end{equation}
Aunque se desprecie el efecto del subsistema térmico en la dinámica del modelo LTI, es interesante seguir estudiando la temperatura del estator en la simulación, por lo que también se considera la siguiente ecuación
\[\dfrac{dT_s(t)}{dt} = \dfrac{1}{C_{ts}}\Big[\frac{3}{2}R_{sREF}\cdot i_{qs}^2(t)-\frac{1}{R_{ts-amb}}\big(T_s(t) - T_{amb}(t)\big)\Big]\]
Definiendo el vector de estados nominal $\mathbf{x}_{LTI} = [\theta_m, \omega_m, i_{qs}]^T$, la representación matricial de estado (separando explícitamente las perturbaciones de carga y gravedad) resulta:

\begingroup
\renewcommand{\arraystretch}{1.8} 
\begin{equation}
\begin{bmatrix}
\dot{\theta}_m (t)\\
\dot{\omega}_m(t)\\
\dot{i}_{qs}(t)
\end{bmatrix}
= 
\underbrace{
\begin{bmatrix}
0 & 1 & 0 \\
0  & -\frac{b_{eq}}{J_{eq}} & \frac{3 P_p \lambda'_m}{2 J_{eq}}     \\
 0 &  -\frac{P_p \lambda'_m}{L_q} & -\frac{R_{s0}}{L_q}
\end{bmatrix}
}_{\mathbf{A}}
\begin{bmatrix}
\theta_m(t) \\
\omega_m (t)\\
i_{qs}(t) 
\end{bmatrix}
+
\underbrace{
\begin{bmatrix}
0  \\
0  \\
\frac{1}{L_q} 
\end{bmatrix}
}_{\mathbf{B_c}}
v_{qs}(t) 
+
\underbrace{
\begin{bmatrix}
0 \\
-\frac{1}{rJ_{eq}} \\
0
\end{bmatrix}
}_{\mathbf{B_d^1}}
T_{ld}(t)
+
\underbrace{
\begin{bmatrix}
    0 \\
    -\frac{g\cdot k_l}{rJ_{eq}} \\
    0
\end{bmatrix}
}_{\mathbf{B_d^2}}
\sin{\left(\frac{\theta_m(t)}{r}\right)}
\end{equation}
\endgroup

Y la ecuación de salida midiendo la posición es:
\begin{equation}
\mathbf{y}(t) = \theta_m(t) = 
\underbrace{
\begin{bmatrix} 1 & 0 & 0 \end{bmatrix}
}_{\mathbf{C}_{LTI}} \mathbf{x}_{LTI} 
+ 
\underbrace{
\begin{bmatrix} 0 & 0 \end{bmatrix}
}_{\mathbf{D}_{LTI}} \mathbf{u}_{LTI}
\end{equation}



% Aquí insertas la imagen de tu diagrama de bloques del sistema LTI.
\begin{figure}[H]
    \centering
    \includegraphics[width=\textwidth]{imagenes/LTI equivalente_nominal.png}
    \caption{Diagrama de bloques de estado del modelo LTI equivalente nominal.}
    \label{fig:diagrama_LTI}
\end{figure}


\subsubsection{Dinámica residual equivalente y Modelo LTI Aumentado}
El modelo nominal desarrollado asume $i_{ds}^r(0) = 0$. En el caso general donde no se cumple esta hipótesis ($i_{ds}^r(0) \neq 0$), la inyección de la ley de control mínima \eqref{eq:vds_minima} en la ecuación del eje directo reduce la misma a una EDO autónoma de primer orden:
\begin{equation}
    \frac{di_{ds}^r(t)}{dt} = -\frac{R_{s0}}{L_d} i_{ds}^r(t) \implies i_{ds}^r(t) = i_{ds}^r(0) e^{-\frac{R_{s0}}{L_d} t}
\end{equation}

Esto demuestra una \textbf{dinámica residual} donde la corriente decae exponencialmente hacia cero. Como el acoplamiento cruzado en el eje $q$ ya fue cancelado matemáticamente por la ley \eqref{eq:vqs_comp}, esta dinámica transitoria ocurre de forma totalmente independiente del resto de los estados. 
Aun asi, esta volveria a acoplar temporalmente la corriente en eje directo con la velocidad, ya que:
\begin{equation*}
    v_{qs}^r(t) = L_q \dfrac{di_{qs}^r(t)}{dt} + R_s(t)i_{qs}^r(t)+\lambda_m'\omega_r(t)+  \underbrace{P_pL_di_{ds}^r(t)\omega_m(t)}_{\text{Acoplamiento}}
\end{equation*}
Para representar el sistema completo, se incorpora esta EDO al modelo LTI, definiendo el vector de estado aumentado $\mathbf{x}_{aug} = [\theta_m, \omega_m, i_{qs}, i_{ds}^r]^T$. 

\subsubsection{Ley de Control Complementaria en el eje $q$}
Si consideramos la dinámica residual analizada en el punto anterior, se puede ver que se producirá acoplamiento cruzado no lineal temporal en el eje de cuadratura debido al término $\omega_r L_d i_{ds}^r$. 

Para eliminar completamente esta dependencia y hacer que el eje $q$ sea lineal, se inyecta una compensación dinámica o \textbf{Ley de control complementaria}. La tensión virtual comandada será:
\begin{equation} \label{eq:vqs_comp}
    v_{qs}^{r*}(t) = v_{qs}(t) + \underbrace{\omega_r(t) L_d i_{ds}^r(t)}_{\text{Cancelación cruzada}}
\end{equation}
Donde $v_{qs}(t)$ pasa a ser la nueva entrada de control lineal del sistema.

El \textbf{Modelo LTI Aumentado} de cuarto orden resultante es:
\begin{equation} \label{eq:sis_escalar_LTI}
\begin{cases}
    \vspace{0.2cm}
    \dfrac{d\theta_m(t)}{dt} = \omega_m(t) \\
    \vspace{0.2cm}
    \dfrac{d\omega_m(t)}{dt} = \dfrac{1}{J_{eq}}\Big[\frac{3 P_p \lambda'_m}{2} i_{qs}(t) - b_{eq} \omega_m(t) - \frac{1}{r} T_{ld}(t) - \frac{g\cdot k_l}{r} \sin\left(\frac{\theta_m(t)}{r}\right)\Big] \\
    \dfrac{di_{qs}(t)}{dt} = \dfrac{1}{L_q}\Big[-P_p \lambda'_m \omega_m(t) - R_{sREF}\cdot i_{qs}(t) + v_{qs}(t)\Big] \\
    \dfrac{di_{ds}(t)}{dt} = \dfrac{-R_{sREF}}{L_d}i_{ds}(t)\\
    \dfrac{dT_s(t)}{dt} = \dfrac{1}{C_{ts}}\Big[\frac{3}{2}R_{sREF}\cdot \big(i_{qs}^2(t)+i_{ds}^2(t)\big)-\frac{1}{R_{ts-amb}}\big(T_s(t) - T_{amb}(t)\big)\Big]
\end{cases}
\end{equation}
\begingroup
\renewcommand{\arraystretch}{1.8} 
\begin{equation} \footnotesize
\begin{bmatrix}
\dot{\theta}_m (t)\\
\dot{\omega}_m(t)\\
\dot{i}_{qs}(t)\\
\dot{i}_{ds}^r(t)
\end{bmatrix}
= 
\begin{bmatrix}
0 & 1 & 0 & 0 \\
0  & -\frac{b_{eq}}{J_{eq}} & \frac{3 P_p \lambda'_m}{2 J_{eq}} & 0 \\
 0 &  -\frac{P_p \lambda'_m}{L_q} & -\frac{R_{s0}}{L_q} & 0 \\
 0 & 0 & 0 & -\frac{R_{s0}}{L_d}
\end{bmatrix}
\begin{bmatrix}
\theta_m(t) \\
\omega_m (t)\\
i_{qs}(t) \\
i_{ds}^r(t)
\end{bmatrix}
+
\begin{bmatrix}
0  \\
0  \\
\frac{1}{L_q} \\
0
\end{bmatrix}
v_{qs}(t) 
+
\begin{bmatrix}
0 \\
-\frac{1}{rJ_{eq}} \\
0 \\
0
\end{bmatrix}
T_{ld}(t)
+
\begin{bmatrix}
    0 \\
    -\frac{k_l}{rJ_{eq}} \\
    0 \\
    0
\end{bmatrix}
\sin{\left(\frac{\theta_m(t)}{r}\right)}
\end{equation}
\endgroup
\subsubsection{Diagrama de bloques de estado (forma desagregada)}
A partir de las ecuaciones matriciales, se construye el diagrama de bloques del sistema linealizado equivalente.
\begin{figure}[H]
    \centering
    \includegraphics[width=\textwidth]{imagenes/LTI equivalente_aumentado.png}
    \caption{Diagrama de bloques de estado del modelo LTI equivalente aumentado.}
    \label{fig:diagrama_LTI_aumentado}
\end{figure}
\subsubsection{Implementación en el modelo global NL}
La implementación física de esta Linealización por Retroalimentación requiere un \textit{controlador parcial} acoplado a la planta NL original. 
\begin{enumerate}
    \item Recibe la acción de control lineal generada por un controlador de movimiento ($v_{qs}$).
    \item Calcula los voltajes de desacoplamiento interno en tiempo real utilizando las lecturas de los sensores idealizados ($\mathbf{i}_{abcs} \rightarrow \mathbf{i}_{qd0s}^r$ y $\omega_m$), aplicando las Ecuaciones \eqref{eq:vds_minima} y \eqref{eq:vqs_comp}.
    \item Estas señales resultantes ($v_{qs}^{r*}$ y $v_{ds}^{r*}$) pasan por un bloque de Transformación de Park Inversa ($\theta_r = P_p \theta_m$).
    \item Finalmente, se envían los voltajes calculados $\mathbf{v}_{abcs}^*$ al inversor, el cual excita físicamente la planta NL.
\end{enumerate}
\begin{figure}[H]
    \centering
    \includegraphics[width=\textwidth]{imagenes/NL con realimentacion NL.png }
    \caption{Diagrama de bloques de estado del modelo LTI equivalente aumentado.}
    \label{fig:diagrama_LTI_aumentado}
\end{figure}
\subsection{Comparación del modelo LPV con el modelo LTI aumentado}
\subsubsection{Comparación: LTI Aumentado vs. LPV con \(I_{ds0}^r \equiv 0\)}

Si bien ambos modelos buscan representar la dinámica del sistema mediante ecuaciones de estado lineales, difieren en su desarrollo y en cómo manejan las restricciones del sistema físico. Las principales diferencias y similitudes se detallan a continuación:

\begin{enumerate}
    \item \textbf{Naturaleza de la linealización:}
    \begin{itemize}
        \item \textit{Modelo LPV:} Se obtiene mediante una expansión en Series de Taylor truncada a primer orden. Es una aproximación de ``pequeña señal'' \cite{franklin_feedback}, válida únicamente en un entorno cercano al punto de equilibrio $(x_0, u_0)$.
        \item \textit{Modelo LTI Aumentado:} Se obtiene mediante Linealización por Retroalimentación NL. Es un modelo algebraicamente exacto sin aproximaciones de Taylor, logrado inyectando tensiones virtuales calculadas dinámicamente para cancelar las no linealidades de la planta.
    \end{itemize}
    
    \item \textbf{Tratamiento de la carga gravitacional (No linealidad mecánica):}
    \begin{itemize}
        \item \textit{Modelo LPV:} El torque gravitacional se linealiza, introduciendo un término de ``rigidez dinámica'' en la matriz de estado del sistema ($A_{2,1} \propto \cos(\theta_0/r)$). Esto provoca que los polos mecánicos cambien drásticamente según la posición geométrica del brazo.
        \item \textit{Modelo LTI Aumentado:} El término gravitacional completo $\sin(\theta_m/r)$ se extrae de la matriz $\mathbf{A}$ y se modela como una perturbación exógena que ingresa al sistema a través de la matriz $\mathbf{B_d^2}$. En consecuencia, la matriz de estado del LTI nominal ($A_{2,1} = 0$) permanece invariante frente a la posición.
    \end{itemize}

    \item \textbf{Desacoplamiento del Subsistema Térmico:}
    \begin{itemize}
        \item \textit{Modelo LPV:} Conserva la influencia de la temperatura como un parámetro de operación. La resistencia $R_{s0}$ dependerá del punto de equilibrio térmico $T_{s0}$ seleccionado para el cálculo del Jacobiano.
        \item \textit{Modelo LTI Aumentado:} Asume un desacoplamiento térmico completo, considerando a $R_s$ como una constante rígida, independizando al subsistema electromecánico de la evolución de la temperatura durante los transitorios rápidos.
    \end{itemize}

    \item \textbf{Similitud estructural al forzar $I_{ds0}^r \equiv 0$:} \\
    A pesar de las diferencias conceptuales, si evaluamos la matriz Jacobiana del modelo LPV forzando el punto de operación $I_{ds0}^r = 0$, observamos una gran convergencia hacia la estructura del modelo LTI Aumentado. 
    
    En el LPV original, los acoplamientos cruzados en la ecuación de tensión y torque generan términos en la matriz dependientes del producto cruzado de los estados. Específicamente, en la derivada de la fuerza contraelectromotriz respecto a la velocidad aparece el término $L_d I_{ds0}^r$. Si imponemos que $I_{ds0}^r = 0$:
    \begin{itemize}
        \item El acoplamiento cruzado de la velocidad sobre el eje $q$ se vuelve puramente proporcional al flujo de los imanes ($\lambda'_m$), coincidiendo exactamente con el término $A_{3,2}$ del LTI Aumentado.
        \item La contribución del torque de reluctancia desaparece en la linealización, quedando el torque dependiendo de forma puramente lineal de $I_{qs0}^r$, al igual que en el modelo LTI.
    \end{itemize}
    
    Por lo que, podemos decir que el Modelo LTI Aumentado puede interpretarse estructuralmente como el caso ideal de un modelo LPV forzado perpetuamente a operar sobre el punto $I_{ds0}^r = 0$, habiendo extraído las variaciones gravitacionales como perturbaciones de entrada.
\end{enumerate}
\subsubsection{Estudio del comportamiento del sistema sometido a variaciones de $i_{ds}^r$}
Al liberar la restricción $I_{ds0}^r = 0$, la corriente del eje directo pasa a ser un parámetro de sintonía que modifica tanto la capacidad estática de la máquina como sus propiedades dinámicas locales (autovalores del Jacobiano). El análisis de esta variación se divide en un enfoque físico-estático y uno dinámico.

\paragraph{1. Análisis Físico: Reforzamiento y Debilitamiento de Campo} \mbox{}\\
La inyección de una corriente continua en el eje directo altera el flujo magnético total de la máquina. El impacto directo se observa en la ecuación del par electromagnético en estado estacionario:
\begin{equation} \label{eq:torque_ids}
    T_m = \frac{3}{2} P_p \Big[ \lambda'_m + (L_d - L_q) I_{ds0}^r \Big] I_{qs0}^r
\end{equation}
Dependiendo del signo de la corriente inyectada, se distinguen dos regímenes de operación:
\begin{itemize}
    \item \textbf{Reforzamiento de campo ($I_{ds0}^r > 0$):} Para motores con polos salientes donde $L_d > L_q$, una corriente positiva genera un torque de reluctancia que se suma al torque de los imanes permanentes. Esto aumenta el torque total por amperio inyectado. Sin embargo, al reforzarse el flujo, la fuerza contraelectromotriz aumenta, lo que reduce la velocidad máxima alcanzable antes de saturar el inversor.\cite{krause_machinery}
    \item \textbf{Debilitamiento de campo ($I_{ds0}^r < 0$):} Una corriente negativa se opone al campo magnético principal de los imanes ($\lambda'_m$). Si bien esto disminuye la capacidad de producción de torque para una misma $I_{qs0}^r$, tiene el efecto beneficioso de reducir drásticamente la fuerza contraelectromotriz. Esto permite que el motor alcance velocidades de giro mucho mayores operando dentro del límite de tensión del inversor.\cite{krause_machinery}
\end{itemize}
\begin{figure}
    \centering
    \includegraphics[width=0.5\linewidth]{imagenes/polosaliente_cilindrico.png}
    \caption{Diagrama representativo de un rotor genérico de polos no-saliente (izq.) y un rotor de polos salientes (der.)\cite{estimation_excitation}}
    \label{fig:placeholder}
\end{figure}

\paragraph{2. Migración de Propiedades Dinámicas (Análisis de Autovalores)} \mbox{}\\
Para evaluar el impacto dinámico, se analiza la migración de los polos del sistema LPV (los autovalores de la matriz $\mathbf{A}_0$). Se fija un punto de operación nominal de velocidad ($\omega_{m0}$) y torque ($I_{qs0}^r$), y se realiza un barrido del parámetro $I_{ds0}^r$ desde la zona de debilitamiento hasta la zona de reforzamiento de campo. 

Sustituyendo estos valores en la matriz Jacobiana $\mathbf{A}_0$ analítica deducida previamente, se obtienen los autovalores para cada caso. Los resultados numéricos más relevantes se resumen en el Cuadro \ref{tab:migracion_polos}.

\begin{table}[H]
\centering

\caption{Migración de los autovalores del modelo LPV ante la variación de $I_{ds0}^r$}
\label{tab:migracion_polos}
\begin{tabular}{|c|c|c|c|c|}
\end{tabular}
\end{table}

A partir de la evaluación numérica de la matriz Jacobiana, se desprenden las siguientes conclusiones sobre la migración de propiedades:
\begin{enumerate}
    \item \textbf{Frecuencia de oscilación:} La variación de $I_{ds0}^r$ afecta drásticamente a los términos cruzados electromecánicos en la matriz $\mathbf{A}_0$. Esto se traduce en una variación significativa de la parte imaginaria de los polos dominantes complejos conjugados. A medida que se pasa de debilitamiento a reforzamiento (aumenta $I_{ds0}^r$), la frecuencia natural de oscilación del sistema aumenta.
    \item \textbf{Estabilidad global:} La parte real de los polos dominantes complejos conjugados (que dicta el amortiguamiento absoluto) se mantiene prácticamente inalterada ante los cambios de $I_{ds0}^r$. Esto demuestra que inyectar corriente en el eje directo no compromete severamente la estabilidad a lazo abierto del sistema.
    \item \textbf{Dinámica rápida:} Los polos puramente reales asociados a la dinámica eléctrica rápida y a la dinámica térmica lenta permanecen virtualmente invariantes frente a los cambios en el punto de operación de $I_{ds0}^r$.
\end{enumerate}

% NOTA PARA TI: Puedes incluir aquí la gráfica (mapa de polos y ceros) generada por MATLAB superponiendo los polos para distintos I_ds0 (como la Figura 14 de tu referencia).
\begin{figure}[H]
    \centering
    % \includegraphics[width=0.8\textwidth]{ruta/migracion_polos.png}
    \caption{Mapa de migración de polos del modelo LPV variando $I_{ds0}^r$.}
    \label{fig:migracion_polos}
\end{figure}
\subsection{Funciones de Transferencia del Modelo LTI Aumentado}
A partir del modelo LTI equivalente aumentado de cuarto orden, se obtienen las funciones de transferencia desde ambas entradas ---la acción de control $v_{qs}(t)$ y la perturbación de carga $T_{ld}(t)$--- hacia la salida de posición $\theta_m(t)$.

\subsubsection{Estado inicial considerado}
El modelo LTI aumentado es un modelo \textbf{incremental} alrededor de un punto de equilibrio. El estado inicial considerado es:
\begin{equation}
    \mathbf{x}_{aug}(0) = \begin{bmatrix} \theta_m(0) \\ \omega_m(0) \\ i_{qs}(0) \\ i_{ds}^r(0) \end{bmatrix} = \begin{bmatrix} 0 \\ 0 \\ 0 \\ 0 \end{bmatrix}
\end{equation}
Además, dado que la gravedad fue extraída de la matriz de estado y se asume compensada mediante \textit{feedforward}, el término $\mathbf{B_d^2}\sin(\theta_m/r)$ no participa en la obtención de las funciones de transferencia lineales. La resistencia del estator se considera constante ($R_s = R_{s0}$) por el desacoplamiento térmico.

\subsubsection{Obtención de las funciones de transferencia}
Las funciones de transferencia se calculan mediante la expresión general:
\begin{equation}
    G(s) = \mathbf{C}(s\mathbf{I} - \mathbf{A})^{-1}\mathbf{B}
\end{equation}

Dado que el estado $i_{ds}^r$ está completamente desacoplado (como se demostrará en la sección siguiente), es suficiente trabajar con el subsistema nominal de tercer orden. Aplicando la transformada de Laplace (con condiciones iniciales nulas) a las ecuaciones escalares del modelo nominal:
\begin{equation}
\begin{cases}
    s\,\Theta_m(s) = \Omega_m(s) \\[6pt]
    s\,\Omega_m(s) = \dfrac{1}{J_{eq}}\left[\dfrac{3 P_p \lambda'_m}{2} I_{qs}(s) - b_{eq}\,\Omega_m(s) - \dfrac{1}{r}\,T_{ld}(s)\right] \\[6pt]
    s\,I_{qs}(s) = \dfrac{1}{L_q}\left[-P_p \lambda'_m\,\Omega_m(s) - R_{s0}\,I_{qs}(s) + V_{qs}(s)\right]
\end{cases}
\end{equation}

De la tercera ecuación se despeja $I_{qs}(s)$:
\begin{equation}\label{eq:iqs_despejada}
    I_{qs}(s) = \frac{V_{qs}(s) - P_p \lambda'_m\,\Omega_m(s)}{L_q\,s + R_{s0}}
\end{equation}

Sustituyendo \eqref{eq:iqs_despejada} en la segunda ecuación y agrupando:
\begin{equation}
    \left(J_{eq}\,s + b_{eq}\right)\Omega_m(s) = \frac{3 P_p \lambda'_m}{2} \cdot \frac{V_{qs}(s) - P_p \lambda'_m\,\Omega_m(s)}{L_q\,s + R_{s0}} - \frac{1}{r}\,T_{ld}(s)
\end{equation}

Multiplicando ambos miembros por $(L_q\,s + R_{s0})$ y reordenando:
\begin{equation}
    \left[(J_{eq}\,s + b_{eq})(L_q\,s + R_{s0}) + \frac{3 P_p^2 \lambda_m'^2}{2}\right]\Omega_m(s) = \frac{3 P_p \lambda'_m}{2}\,V_{qs}(s) - \frac{L_q\,s + R_{s0}}{r}\,T_{ld}(s)
\end{equation}

Finalmente, dado que $\Theta_m(s) = \Omega_m(s)/s$, y expandiendo el término entre corchetes se obtiene el denominador común:
\begin{equation}
    \Delta(s) = s\left[J_{eq} L_q\,s^2 + (b_{eq} L_q + J_{eq} R_{s0})\,s + b_{eq} R_{s0} + \frac{3 P_p^2 \lambda_m'^2}{2}\right]
\end{equation}

donde el factor $s$ proviene de la integración $\Theta_m = \Omega_m / s$. Dividiendo numerador y denominador por $J_{eq} L_q$ se obtienen las expresiones finales.

\vspace{0.3cm}

El polinomio característico del modelo aumentado es:
\begin{equation}
    \det(s\mathbf{I} - \mathbf{A}_{aug}) = \underbrace{\left(s + \frac{R_{s0}}{L_d}\right)}_{\text{Polo de } i_{ds}^r} \cdot \underbrace{\left[ s^3 + \left(\frac{b_{eq}}{J_{eq}} + \frac{R_{s0}}{L_q}\right)s^2 + \left(\frac{b_{eq} R_{s0}}{J_{eq} L_q} + \frac{3 P_p^2 \lambda_m'^2}{2 J_{eq} L_q}\right)s \right]}_{\Delta(s) \text{ --- Polinomio del subsistema nominal}}
\end{equation}

Sin embargo, al evalar la resolvente $(s\mathbf{I}-\mathbf{A})^{-1}$ y proyectar sobre las columnas de entrada $\mathbf{B_c}$ y $\mathbf{B_d}$, el factor $(s + R_{s0}/L_d)$ se \textbf{cancela} en ambas funciones de transferencia. Esto se debe a que el estado $i_{ds}^r$ no es excitado por ninguna de las dos entradas ($B_c(4) = B_d(4) = 0$) y no es observado por la salida ($C(4) = 0$). Por tanto, las funciones de transferencia resultan de \textbf{orden 3}, con denominador $\Delta(s)$:

\vspace{0.3cm}
\textbf{Función de transferencia de control} ($v_{qs} \to \theta_m$):
\begin{equation} \label{eq:tf_vqs}
\boxed{
    G_1(s) = \frac{\theta_m(s)}{v_{qs}(s)} = \frac{\dfrac{3 P_p \lambda'_m}{2 J_{eq} L_q}}{s^3 + \left(\dfrac{b_{eq}}{J_{eq}} + \dfrac{R_{s0}}{L_q}\right)s^2 + \left(\dfrac{b_{eq} R_{s0}}{J_{eq} L_q} + \dfrac{3 P_p^2 \lambda_m'^2}{2 J_{eq} L_q}\right)s}
}
\end{equation}

\vspace{0.3cm}
\textbf{Función de transferencia de perturbación} ($T_{ld} \to \theta_m$):
\begin{equation} \label{eq:tf_Tld}
\boxed{
    G_2(s) = \frac{\theta_m(s)}{T_{ld}(s)} = \frac{-\dfrac{1}{r J_{eq}}\left(s + \dfrac{R_{s0}}{L_q}\right)}{s^3 + \left(\dfrac{b_{eq}}{J_{eq}} + \dfrac{R_{s0}}{L_q}\right)s^2 + \left(\dfrac{b_{eq} R_{s0}}{J_{eq} L_q} + \dfrac{3 P_p^2 \lambda_m'^2}{2 J_{eq} L_q}\right)s}
}
\end{equation}

Ambas funciones de transferencia comparten el mismo polinomio característico $\Delta(s)$, el cual presenta un polo en el origen ($s=0$) correspondiente a la integración inherente de velocidad a posición. Esto indica que el sistema a lazo abierto es de \textbf{tipo 1}: sigue rampas de posición con error estacionario finito, pero no alcanza error nulo ante entrada escalón sin realimentación.

Notar que $G_2(s)$ posee un \textbf{cero} en $s = -R_{s0}/L_q$, que coincide con el polo eléctrico del eje $q$. Esto se interpreta físicamente como: ante una perturbación de torque de carga, la corriente $i_{qs}$ reacciona inicialmente absorbiendo parte del impacto, atenuando la transmisión del disturbio hacia la posición en frecuencias altas.

\subsubsection{Estado interno no reflejado en las funciones de transferencia}
El estado $i_{ds}^r$ (cuarto estado del modelo aumentado) \textbf{no aporta ni se ve reflejado} en ninguna de las funciones de transferencia obtenidas. Formalmente, esto se verifica analizando las propiedades estructurales del sistema:

\begin{enumerate}
    \item \textbf{No controlabilidad:} Ninguna de las entradas ($v_{qs}$, $T_{ld}$) excita al estado $i_{ds}^r$, ya que las componentes correspondientes de las matrices de entrada son nulas: $B_c(4) = 0$ y $B_d(4) = 0$.
    \item \textbf{No observabilidad:} La salida $\theta_m$ no depende de $i_{ds}^r$, ya que $C(4) = 0$.
    \item \textbf{Desacoplamiento dinámico:} La cuarta fila de $\mathbf{A}$ solo contiene el término $-R_{s0}/L_d$ en la diagonal, sin acoplamientos con los demás estados. La dinámica de $i_{ds}^r$ es completamente autónoma.
\end{enumerate}

Consecuentemente, el polo $\lambda_4 = -R_{s0}/L_d$ del sistema en espacio de estados se cancela algebraicamente al convertir a función de transferencia. La dinámica de $i_{ds}^r$ es puramente una \textbf{respuesta libre}:
\begin{equation}
    i_{ds}^r(t) = i_{ds}^r(0) \, e^{-\frac{R_{s0}}{L_d} t}
\end{equation}
que solo se manifiesta si la condición inicial es no nula, decayendo exponencialmente con constante de tiempo $\tau_d = L_d / R_{s0}$. Esto confirma que el \textbf{modelo nominal de tercer orden} $\mathbf{x}_{LTI} = [\theta_m, \omega_m, i_{qs}]^T$ es suficiente para caracterizar completamente la relación entrada-salida del sistema.

\subsection{Análisis de Estabilidad a Lazo Abierto}

\subsubsection{Autovalores del modelo aumentado}
Los autovalores del sistema se obtienen a partir del polinomio característico de la matriz $\mathbf{A}_{aug}$ (4$\times$4). Como se demostró en la sección anterior, este se factoriza en dos términos independientes:
\begin{equation}
    \det(s\mathbf{I} - \mathbf{A}_{aug}) = \underbrace{\left(s + \frac{R_{s0}}{L_d}\right)}_{\lambda_4} \cdot \underbrace{s}_{\lambda_1} \cdot \underbrace{\left(s^2 + a_2\,s + a_1\right)}_{\lambda_{2,3}}
\end{equation}
donde los coeficientes del polinomio cuadrático son:
\begin{equation}
    a_2 = \frac{b_{eq}}{J_{eq}} + \frac{R_{s0}}{L_q}, \qquad a_1 = \frac{b_{eq} R_{s0}}{J_{eq} L_q} + \frac{3 P_p^2 \lambda_m'^2}{2 J_{eq} L_q}
\end{equation}

Evaluando numéricamente con los parámetros nominales ($m_l = 0$, $J_{l,nom} = 0{,}0833$ kg$\cdot$m$^2$):

\begin{table}[H]
\centering
\caption{Autovalores del modelo LTI aumentado a lazo abierto (parámetros nominales).}
\label{tab:autovalores_LA}
\renewcommand{\arraystretch}{1.4}
\begin{tabular}{|c|c|c|c|c|}
\hline
\textbf{Autovalor} & \textbf{Valor} & \textbf{Subsistema} & $\boldsymbol{\omega_n}$ \textbf{[rad/s]} & $\boldsymbol{\zeta}$ \\
\hline
$\lambda_1$ & $0$ & Integrador (posición) & --- & --- \\
\hline
$\lambda_{2,3}$ & $-88{,}5 \pm j\,150{,}0$ & Electromecánico & $174{,}1$ & $0{,}508$ \\
\hline
$\lambda_4$ & $-154{,}5$ & Residual $i_{ds}^r$ & --- & $\tau_d = 6{,}5$ ms \\
\hline
\end{tabular}
\end{table}

\paragraph{Identificación física de cada autovalor:} \mbox{}

\begin{itemize}
    \item \textbf{$\lambda_1 = 0$ (Polo en el origen):} Corresponde a la relación cinemática $\dot{\theta}_m = \omega_m$. La posición es la integral de la velocidad, por lo que el sistema posee un \textbf{integrador puro}. Esto implica que ante una entrada constante, la posición crece indefinidamente sin realimentación.

    \item \textbf{$\lambda_{2,3} = -88{,}5 \pm j\,150{,}0$ (Par complejo conjugado):} Representan el \textbf{modo electromecánico acoplado} entre la corriente $i_{qs}$ y la velocidad $\omega_m$. La parte imaginaria indica un comportamiento oscilatorio subamortiguado ($\zeta \approx 0{,}51$), producto del intercambio de energía entre el campo magnético ($L_q$, $\lambda'_m$) y la inercia mecánica ($J_{eq}$). La frecuencia de oscilación amortiguada es $\omega_d = 150{,}0$ rad/s ($\approx 23{,}9$ Hz), y la envolvente decae con constante de tiempo $\tau_{em} = 1/88{,}5 \approx 11{,}3$ ms.

    \item \textbf{$\lambda_4 = -R_{s0}/L_d = -154{,}5$ rad/s (Polo real negativo):} Corresponde a la dinámica residual autónoma de $i_{ds}^r$. Como se demostró, este polo \textbf{no aparece} en las funciones de transferencia. Su constante de tiempo es $\tau_d = L_d/R_{s0} \approx 6{,}5$ ms.
\end{itemize}

\subsubsection{Ceros de las funciones de transferencia}

\begin{itemize}
    \item \textbf{$G_1(s) = \theta_m / v_{qs}$:} No posee ceros finitos. Es un sistema de \textbf{fase mínima} puro de tipo 1 con exceso polar de 3. La respuesta ante un escalón de tensión será una rampa modulada por la dinámica oscilatoria de los polos electromecánicos.

    \item \textbf{$G_2(s) = \theta_m / T_{ld}$:} Posee un cero real negativo en:
    \begin{equation}
        z_1 = -\frac{R_{s0}}{L_q} = -175{,}9 \text{ rad/s} \qquad (\tau_z = 5{,}7 \text{ ms})
    \end{equation}
    Este cero se ubica muy cerca del polo electromecánico dominante. Físicamente, representa que el subsistema eléctrico ($L_q$, $R_{s0}$) ``absorbe'' parcialmente las perturbaciones de torque antes de que estas se transmitan completamente a la posición, actuando como un filtro natural de alta frecuencia sobre la perturbación.
\end{itemize}

\subsubsection{Estabilidad}

\paragraph{Estabilidad interna (Lyapunov):} \mbox{} \\
El modelo LTI aumentado posee un autovalor en $s = 0$ (simple) y los restantes con parte real estrictamente negativa. Por tanto, el sistema es \textbf{marginalmente estable} (o estable en el sentido de Lyapunov, pero \textbf{no} asintóticamente estable). Una perturbación en la posición no se atenúa: el sistema ``recuerda'' cualquier desplazamiento permanentemente.

\paragraph{Estabilidad BIBO (entrada-salida acotada):} \mbox{} \\
Para la salida $\theta_m$:
\begin{itemize}
    \item Desde $v_{qs}$: \textbf{No es BIBO estable}. El polo en el origen de $G_1(s)$ implica que un escalón constante de tensión produce una posición que crece sin cota (rampa).
    \item Desde $T_{ld}$: \textbf{No es BIBO estable}. Análogamente, un torque de perturbación constante genera una deriva indefinida de la posición.
\end{itemize}

\paragraph{Estabilidad parcial:} \mbox{} \\
Si se analiza únicamente el subsistema de velocidad y corriente ($\omega_m$, $i_{qs}$), los polos electromecánicos $\lambda_{2,3}$ tienen parte real negativa, garantizando que las \textbf{variables de velocidad y corriente convergen} a valores estacionarios finitos ante entradas acotadas. Solo la posición es la variable inestable (integradora).

Esta estructura es característica de los sistemas de \textbf{regulación de posición}: se requiere necesariamente un lazo cerrado de realimentación para estabilizar la posición. El controlador PID que se diseñará en la sección 5.2 cumplirá exactamente esta función.

\begin{figure}[H]
    \centering
    \includegraphics[width=\textwidth]{imagenes/polos_ceros_LA.png}
    \caption{Mapa de polos y ceros a lazo abierto para ambas funciones de transferencia del modelo LTI aumentado. Se indica con $\times$ roja el polo cancelado de $i_{ds}^r$.}
    \label{fig:polos_ceros_LA}
\end{figure}

\subsubsection{Migración de propiedades ante variación de $R_s$}
La resistencia del estator varía con la temperatura según $R_s(T_s) = R_{s0}(1 + \alpha_{Cu}(T_s - T_{s,ref}))$. Para evaluar su efecto sobre la dinámica, se realiza un barrido en el rango $[0{,}8\,R_{s0},\; 1{,}2\,R_{s0}]$.

\begin{table}[H]
\centering
\caption{Migración de polos y ceros ante variación de $R_s$.}
\label{tab:migracion_Rs}
\renewcommand{\arraystretch}{1.2}
\footnotesize
\begin{tabular}{|c|c|c|c|c|}
\hline
$R_s$ [$\Omega$] & Polos $\lambda_{2,3}$ & Cero $G_2$ & $\omega_n$ [rad/s] & $\zeta$ \\
\hline
0.816 & $-70{,}9 \pm j\,158{,}9$ & $-140{,}69$ & $174{,}0$ & $0{,}407$ \\
0.850 & $-73{,}8 \pm j\,157{,}6$ & $-146{,}55$ & $174{,}0$ & $0{,}424$ \\
0.884 & $-76{,}8 \pm j\,156{,}2$ & $-152{,}41$ & $174{,}0$ & $0{,}441$ \\
0.918 & $-79{,}7 \pm j\,154{,}7$ & $-158{,}28$ & $174{,}0$ & $0{,}458$ \\
0.952 & $-82{,}6 \pm j\,153{,}2$ & $-164{,}14$ & $174{,}1$ & $0{,}475$ \\
0.986 & $-85{,}6 \pm j\,151{,}6$ & $-170{,}00$ & $174{,}1$ & $0{,}491$ \\
\textbf{1.020} & $\mathbf{-88{,}5 \pm j\,149{,}9}$ & $\mathbf{-175{,}86}$ & $\mathbf{174{,}1}$ & $\mathbf{0{,}508}$ \\
1.054 & $-91{,}4 \pm j\,148{,}2$ & $-181{,}72$ & $174{,}1$ & $0{,}525$ \\
1.088 & $-94{,}3 \pm j\,146{,}4$ & $-187{,}59$ & $174{,}1$ & $0{,}542$ \\
1.122 & $-97{,}3 \pm j\,144{,}5$ & $-193{,}45$ & $174{,}2$ & $0{,}559$ \\
1.156 & $-100{,}2 \pm j\,142{,}5$ & $-199{,}31$ & $174{,}2$ & $0{,}575$ \\
1.190 & $-103{,}1 \pm j\,140{,}4$ & $-205{,}17$ & $174{,}2$ & $0{,}592$ \\
1.224 & $-106{,}1 \pm j\,138{,}2$ & $-211{,}03$ & $174{,}2$ & $0{,}609$ \\
\hline
\end{tabular}
\end{table}

Del análisis se concluye:
\begin{itemize}
    \item La \textbf{frecuencia natural} $\omega_n$ permanece prácticamente constante ($\approx 174$ rad/s), ya que depende principalmente del acoplamiento electromecánico $P_p\lambda'_m$ y la inercia $J_{eq}$, no de $R_s$.
    \item El \textbf{amortiguamiento} $\zeta$ aumenta con $R_s$ (de $0{,}41$ a $0{,}61$), lo que se explica porque una mayor resistencia disipa más rápidamente la energía almacenada en el campo magnético.
    \item El \textbf{cero} de $G_2(s)$ migra proporcionalmente a $R_s$ ($z = -R_s/L_q$), desplazándose hacia la izquierda con el aumento de temperatura.
\end{itemize}

\begin{figure}[H]
    \centering
    \includegraphics[width=\textwidth]{imagenes/migracion_Rs.png}
    \caption{Migración de polos y ceros en el plano $s$ (izquierda) y variación de $\omega_n$ y $\zeta$ (derecha) ante cambios en $R_s$.}
    \label{fig:migracion_Rs}
\end{figure}

\subsubsection{Migración de propiedades ante variación de carga ($m_l$)}
La masa de carga útil $m_l \in [0,\; 1{,}5]$ kg modifica la inercia equivalente $J_{eq} = J_m + J_l/r^2$ y el coeficiente gravitacional $k_l$. Al aumentar $m_l$, la inercia crece y el acoplamiento electromecánico se ``diluye'', modificando los autovalores.

\begin{table}[H]
\centering
\caption{Migración de polos ante variación de masa de carga $m_l$.}
\label{tab:migracion_ml}
\renewcommand{\arraystretch}{1.2}
\footnotesize
\begin{tabular}{|c|c|c|c|c|}
\hline
$m_l$ [kg] & $J_{eq}$ [kg$\cdot$m$^2$] & Polos $\lambda_{2,3}$ & $\omega_n$ [rad/s] & $\zeta$ \\
\hline
\textbf{0.00} & $\mathbf{1{,}98 \times 10^{-5}}$ & $\mathbf{-88{,}5 \pm j\,149{,}9}$ & $\mathbf{174{,}1}$ & $\mathbf{0{,}508}$ \\
0.20 & $2{,}33 \times 10^{-5}$ & $-88{,}4 \pm j\,134{,}1$ & $160{,}6$ & $0{,}551$ \\
0.40 & $2{,}67 \times 10^{-5}$ & $-88{,}3 \pm j\,121{,}0$ & $149{,}8$ & $0{,}590$ \\
0.60 & $3{,}02 \times 10^{-5}$ & $-88{,}3 \pm j\,109{,}8$ & $140{,}9$ & $0{,}627$ \\
0.80 & $3{,}37 \times 10^{-5}$ & $-88{,}3 \pm j\,100{,}1$ & $133{,}5$ & $0{,}661$ \\
1.00 & $3{,}71 \times 10^{-5}$ & $-88{,}2 \pm j\,91{,}4$ & $127{,}1$ & $0{,}694$ \\
1.20 & $4{,}06 \times 10^{-5}$ & $-88{,}2 \pm j\,83{,}6$ & $121{,}5$ & $0{,}726$ \\
1.50 & $4{,}58 \times 10^{-5}$ & $-88{,}2 \pm j\,72{,}9$ & $114{,}4$ & $0{,}771$ \\
\hline
\end{tabular}
\end{table}

\begin{figure}[H]
    \centering
    \includegraphics[width=\textwidth]{imagenes/migracion_ml.png}
    \caption{Migración de polos ante variación de masa de carga $m_l$ (izquierda) y variación de $\omega_n$ y $\zeta$ (derecha).}
    \label{fig:migracion_ml}
\end{figure}

Del análisis se observa que al aumentar $m_l$:
\begin{itemize}
    \item La \textbf{frecuencia natural} $\omega_n$ disminuye, ya que la mayor inercia reduce la relación torque/inercia del sistema.
    \item El \textbf{amortiguamiento} $\zeta$ aumenta, porque la fricción equivalente $b_{eq}$ se mantiene mientras la frecuencia natural baja, incrementando la relación $\zeta = a_2/(2\omega_n)$.
    \item El polo en el origen y el cero de $G_2$ permanecen invariantes frente a $m_l$.
\end{itemize}

Estas variaciones justifican la necesidad de un controlador robusto capaz de mantener el desempeño dentro de todo el rango de incertidumbre de carga.

\subsubsection{Diagramas de Bode a lazo abierto}
\begin{figure}[H]
    \centering
    \includegraphics[width=\textwidth]{imagenes/bode_LA.png}
    \caption{Diagramas de Bode de las funciones de transferencia $G_1(s)$ y $G_2(s)$ a lazo abierto.}
    \label{fig:bode_LA}
\end{figure}

\subsection{Análisis de Observabilidad}
Se analiza si es posible reconstruir el vector de estado completo del modelo LTI aumentado a partir de la medición de una única salida escalar.

\subsubsection{Observabilidad desde $\theta_m(t)$ (sensor de posición)}
La matriz de observabilidad para el sistema $(\mathbf{A}_{aug}, \mathbf{C})$ con $\mathbf{C} = \begin{bmatrix} 1 & 0 & 0 & 0 \end{bmatrix}$ se construye como:
\begin{equation}
    \mathcal{O} = \begin{bmatrix} \mathbf{C} \\ \mathbf{C}\mathbf{A} \\ \mathbf{C}\mathbf{A}^2 \\ \mathbf{C}\mathbf{A}^3 \end{bmatrix}
\end{equation}

Evaluando fila por fila a partir de la estructura de $\mathbf{A}_{aug}$:
\begin{itemize}
    \item $\mathbf{C} = \begin{bmatrix} 1 & 0 & 0 & 0 \end{bmatrix}$ --- se observa $\theta_m$ directamente.
    \item $\mathbf{C}\mathbf{A} = \begin{bmatrix} 0 & 1 & 0 & 0 \end{bmatrix}$ --- se observa $\omega_m$ (derivada de $\theta_m$).
    \item $\mathbf{C}\mathbf{A}^2 = \begin{bmatrix} 0 & -\frac{b_{eq}}{J_{eq}} & \frac{3P_p\lambda'_m}{2J_{eq}} & 0 \end{bmatrix}$ --- combinación de $\omega_m$ e $i_{qs}$.
    \item $\mathbf{C}\mathbf{A}^3$: combinación que involucra $\omega_m$ e $i_{qs}$, pero \textbf{nunca} $i_{ds}^r$.
\end{itemize}

La cuarta columna de $\mathcal{O}$ es idénticamente nula, ya que $i_{ds}^r$ no acopla con ningún estado observable. Por lo tanto:
\begin{equation}
    \text{rango}(\mathcal{O}) = 3 < 4 = n
\end{equation}

\textbf{Conclusión:} El sistema \textbf{no es completamente observable} desde $\theta_m(t)$. El estado no observable es $i_{ds}^r$. Los tres estados restantes ($\theta_m$, $\omega_m$, $i_{qs}$) sí son observables, lo que permite diseñar un observador reducido para estimar $\omega_m$ e $i_{qs}$ a partir de la medición de posición.

Esto es consistente con el análisis de las funciones de transferencia: dado que $i_{ds}^r$ no se manifiesta en la salida $\theta_m$, es imposible reconstruirlo a partir de dicha medición.

\subsubsection{Observabilidad desde $\omega_m(t)$ (tacogenerador)}
Si en lugar de medir la posición se mide la velocidad angular, la matriz de salida es $\mathbf{C}_\omega = \begin{bmatrix} 0 & 1 & 0 & 0 \end{bmatrix}$. Evaluando:
\begin{itemize}
    \item $\mathbf{C}_\omega = \begin{bmatrix} 0 & 1 & 0 & 0 \end{bmatrix}$ --- se observa $\omega_m$.
    \item $\mathbf{C}_\omega\mathbf{A} = \begin{bmatrix} 0 & -\frac{b_{eq}}{J_{eq}} & \frac{3P_p\lambda'_m}{2J_{eq}} & 0 \end{bmatrix}$ --- combinación de $\omega_m$ e $i_{qs}$.
    \item $\mathbf{C}_\omega\mathbf{A}^2$: combinación de $\omega_m$ e $i_{qs}$.
    \item $\mathbf{C}_\omega\mathbf{A}^3$: combinación de $\omega_m$ e $i_{qs}$.
\end{itemize}

Nuevamente, la cuarta columna es nula ($i_{ds}^r$ sigue sin ser observable), pero además la primera columna también es nula: la posición $\theta_m$ nunca aparece en las potencias de $\mathbf{C}_\omega\mathbf{A}^k$ porque $\theta_m$ solo se acopla unidireccionalmente como integrador de $\omega_m$.

\begin{equation}
    \text{rango}(\mathcal{O}_\omega) = 2 < 4 = n
\end{equation}

\textbf{Conclusión:} Midiendo $\omega_m$ se pierden \textbf{dos} estados: $\theta_m$ e $i_{ds}^r$. La posición no es observable desde la velocidad porque la integración ``borra'' la condición inicial ---es decir, infinitas trayectorias de posición producen la misma velocidad. Por tanto, medir $\theta_m$ con un \textit{encoder} es estrictamente más informativo que medir $\omega_m$ con un tacogenerador para este sistema.

\subsection{Análisis de Controlabilidad}
Se analiza si es posible llevar el vector de estado completo a cualquier valor deseado mediante la entrada de control $v_{qs}(t)$, sin considerar la perturbación $T_{ld}(t)$.

\subsubsection{Controlabilidad desde $v_{qs}(t)$}
La matriz de controlabilidad para el sistema $(\mathbf{A}_{aug}, \mathbf{B}_c)$ con $\mathbf{B}_c = \begin{bmatrix} 0 & 0 & 1/L_q & 0 \end{bmatrix}^T$ es:
\begin{equation}
    \mathcal{C} = \begin{bmatrix} \mathbf{B}_c & \mathbf{A}\mathbf{B}_c & \mathbf{A}^2\mathbf{B}_c & \mathbf{A}^3\mathbf{B}_c \end{bmatrix}
\end{equation}

Evaluando columna por columna:
\begin{itemize}
    \item $\mathbf{B}_c = \begin{bmatrix} 0 \\ 0 \\ 1/L_q \\ 0 \end{bmatrix}$ --- la entrada actúa directamente sobre $i_{qs}$.
    \item $\mathbf{A}\mathbf{B}_c = \begin{bmatrix} 0 \\ \frac{3P_p\lambda'_m}{2J_{eq}L_q} \\ -\frac{R_{s0}}{L_q^2} \\ 0 \end{bmatrix}$ --- la acción se propaga a $\omega_m$ a través del torque.
    \item $\mathbf{A}^2\mathbf{B}_c$: la acción alcanza $\theta_m$ (vía la cadena $i_{qs} \to \omega_m \to \theta_m$).
    \item $\mathbf{A}^3\mathbf{B}_c$: combinación de los tres estados del subsistema nominal.
\end{itemize}

La cuarta fila de todas las columnas de $\mathcal{C}$ es idénticamente nula, ya que la entrada $v_{qs}$ nunca llega al estado $i_{ds}^r$ ($B_c(4) = 0$) y ningún otro estado acopla hacia $i_{ds}^r$ (la fila 4 de $\mathbf{A}$ es diagonal). Por lo tanto:
\begin{equation}
    \text{rango}(\mathcal{C}) = 3 < 4 = n
\end{equation}

\textbf{Conclusión:} El sistema \textbf{no es completamente controlable} desde $v_{qs}(t)$. El estado no controlable es $i_{ds}^r$.

\subsubsection{Incorporación de una entrada adicional}
El estado $i_{ds}^r$ podría controlarse si se agrega una segunda entrada manipulable: la tensión del eje directo $v_{ds}^r(t)$. En el modelo aumentado, esto equivale a extender la matriz de entrada:
\begin{equation}
    \mathbf{B}_{ext} = \begin{bmatrix} \mathbf{B}_c & \mathbf{B}_{ds} \end{bmatrix} =
    \begin{bmatrix}
    0 & 0 \\ 0 & 0 \\ 1/L_q & 0 \\ 0 & 1/L_d
    \end{bmatrix}
\end{equation}

Con esta extensión, la matriz de controlabilidad alcanza rango 4 y el sistema pasa a ser completamente controlable. Sin embargo, en la práctica esta entrada $v_{ds}^r$ ya está siendo utilizada por la \textbf{ley de control mínima} \eqref{eq:vds_minima} para forzar $i_{ds}^r \to 0$. Es decir, la propia estrategia de control vectorial ya ``controla'' $i_{ds}^r$ llevándola exponencialmente a cero; simplemente no lo hace a través de la estructura de realimentación de estado convencional, sino mediante la inyección de una tensión de desacoplamiento.

\subsubsection{Resumen de propiedades estructurales}

\begin{table}[H]
\centering
\caption{Resumen de controlabilidad y observabilidad del modelo LTI aumentado.}
\label{tab:ctrl_obs}
\renewcommand{\arraystretch}{1.3}
\begin{tabular}{|l|c|c|c|c|}
\hline
\textbf{Estado} & \textbf{Controlable} ($v_{qs}$) & \textbf{Observable} ($\theta_m$) & \textbf{Observable} ($\omega_m$) & \textbf{En las FTs} \\
\hline
$\theta_m$   & Sí & Sí (directo) & No & Sí \\
$\omega_m$   & Sí & Sí           & Sí (directo) & Sí \\
$i_{qs}$     & Sí (directo) & Sí      & Sí & Sí \\
$i_{ds}^r$   & No & No            & No & No \\
\hline
\end{tabular}
\end{table}

El estado $i_{ds}^r$ es simultáneamente no controlable y no observable, confirmando su naturaleza de \textbf{modo oculto} del sistema. Este resultado es coherente con la cancelación polo-cero identificada en las funciones de transferencia y con el desacoplamiento estructural de la cuarta fila y columna de la matriz $\mathbf{A}_{aug}$.

\bibliographystyle{ieeetr}
\bibliography{references}
\end{document}